% Generated draft: requires mathematical and source review.
\documentclass[11pt]{article}
\usepackage[margin=25mm]{geometry}
\usepackage{fontspec}
\setmainfont{Latin Modern Roman}
\usepackage{amsmath,amssymb,mathtools,unicode-math}
\setmathfont{Latin Modern Math}
\usepackage{enumitem,needspace,xurl,hyperref}
\setlength{\parindent}{0pt}
\setlength{\parskip}{5pt}
\setlength{\emergencystretch}{3em}
\newcommand{\N}{\mathbb{N}}
\newcommand{\Z}{\mathbb{Z}}
\newcommand{\Q}{\mathbb{Q}}
\newcommand{\R}{\mathbb{R}}
\newcommand{\C}{\mathbb{C}}
\newcommand{\F}{\mathbb{F}}
\newcommand{\A}{\mathbb{A}}
\newcommand{\PP}{\mathbb{P}}
\newcommand{\E}{\mathbb{E}}
\newcommand{\eps}{\varepsilon}
\newcommand{\dd}{\,\mathrm{d}}
\begin{document}
\section*{UnsolvedMath ID: 30002929}
\textbf{Source problem code:} catalogue entry
\subsection{Definitions and mathematical statement}
The supplied conjecture is: for superfluids, the total momentum is zero with positive asymptotic probability. The record does not specify a mathematical definition of a superfluid, the probability law of total momentum, or the asymptotic regime, so it does not determine a unique fully quantified statement.

A conditional formulation can be given once these data are supplied. Let \(n\) index a specified sequence of systems in a specified superfluid regime. Let \((\Omega_n,\mathcal F_n,\mathbb P_n)\) describe the total-momentum measurement distribution, and let \(\mathbf P_n:\Omega_n\to\R^d\) denote total momentum, where the spatial dimension \(d\geq 1\) must also be specified. Set
\[
p_n:=\mathbb P_n\bigl(\mathbf P_n=\mathbf 0\bigr).
\]
If ``positive asymptotic probability'' means convergence of the probabilities of exactly zero momentum, the proposed assertion is
\[
\lim_{n\to\infty}p_n=p_0\qquad\text{for some }p_0>0.
\]
This is a conditional statement template, not a determination of the source's intended model or limit. The supplied record does not establish whether it instead intends \(\liminf_{n\to\infty}p_n>0\), or positive mass at zero in a limiting momentum distribution. In particular, if the momentum laws converge weakly to a probability measure \(\mu\), the assertion \(\mu(\{\mathbf 0\})>0\) is generally different from the displayed assertion. The class of systems over which the conjecture is quantified, the equilibrium state or ensemble, interactions, boundary conditions, reference frame, temperature and density conditions, and limiting procedure remain unspecified.
\subsection{Short English statement}
The conjecture asks whether a superfluid has strictly positive asymptotic probability of exactly zero total momentum. A precise statement requires a specified superfluid model, state, reference frame, and limiting regime, together with clarification of whether positivity concerns exact finite-system probabilities or an atom of a limiting distribution. The supplied record labels the problem open; this draft does not independently verify that status.
\subsection{Sources}
\begin{itemize}
\item Supplied local UnsolvedMath registry record, permanent ID \(30002929\), title ``Positive Probability of Zero Total Superfluid Momentum,'' source problem code \texttt{catalogue entry}, in the Oberwolfach Reports Open Problems collection. The record supplies proposed year 2015 and no named proposer.
\item Supplied report citation: \emph{Interplay of Analysis and Probability in Applied Mathematics} (2016), Oberwolfach Reports, pp.\ 1989--2064, DOI \texttt{10.4171/owr/2015/35}. \url{https://doi.org/10.4171/owr/2015/35}. The background describes Oberwolfach Reports as documenting workshops and their open-problem discussions.
\item Supplied literature reference: Suto, \emph{The total momentum of quantum fluids}. \url{https://arxiv.org/abs/1504.06528}. The supplied assessment attributes the zero-total-momentum conjecture for superfluids to Suto; it does not identify Suto as the proposer of this registry entry.
\item Literature assessment checked at \texttt{2026-08-21}: the supplied curation labels the problem open and reports that subsequent literature located through its cutoff does not prove the stated positive-probability claim. This is supplied dated curation, not an independently verified literature assessment. The background expressly requests independent verification of time-sensitive status and cited claims.
\item The registry labels the statement \texttt{corrected\_verified} and says its mathematical content was preserved while notation, typography, or omitted definitions were repaired for self-contained publication. These are supplied editorial assertions; this draft does not independently verify them or the linked references.
\end{itemize}
\end{document}
